\par\medskip{\large\bfseries\color{mauPhieu} Bài 1. Tọa độ của véctơ}\par\smallskip
%%% BAI 0H9?1
\textbf{1. Kiến thức cần nhớ}
\begin{itemize}
\item Hệ trục tọa độ $Oxy$ có hai véctơ đơn vị $\overrightarrow{i}, \overrightarrow{j}$ tương ứng trên $Ox, Oy$.
\item Tọa độ véctơ: $\overrightarrow{u} = (x;y) \Leftrightarrow \overrightarrow{u} = x\overrightarrow{i} + y\overrightarrow{j}$.
\item Tọa độ điểm: $M(x_M; y_M) \Leftrightarrow \overrightarrow{OM} = x_M\overrightarrow{i} + y_M\overrightarrow{j}$.
\item Véctơ giữa hai điểm $A, B$: $\overrightarrow{AB} = (x_B - x_A; y_B - y_A)$.
\item Các phép toán véctơ: Cho $\overrightarrow{a} = (x_1; y_1)$ và $\overrightarrow{b} = (x_2; y_2)$.
  \begin{itemize}
  \item $\overrightarrow{a} \pm \overrightarrow{b} = (x_1 \pm x_2; y_1 \pm y_2)$.
  \item $k\overrightarrow{a} = (kx_1; ky_1)$ với $k \in \mathbb{R}$.
  \item $\overrightarrow{a} = \overrightarrow{b} \Leftrightarrow x_1 = x_2$ và $y_1 = y_2$.
  \end{itemize}
\item Hai véctơ $\overrightarrow{a}, \overrightarrow{b}$ ($\overrightarrow{b} \neq \overrightarrow{0}$) cùng phương $\Leftrightarrow \exists k \in \mathbb{R}$ sao cho $\overrightarrow{a} = k\overrightarrow{b} \Leftrightarrow x_1 = k x_2$ và $y_1 = k y_2$.
\item Tọa độ trung điểm $I$ của $AB$: $\boxed{x_I = \frac{x_A + x_B}{2}; y_I = \frac{y_A + y_B}{2}}$.
\item Tọa độ trọng tâm $G$ của $\Delta ABC$: $\boxed{x_G = \frac{x_A + x_B + x_C}{3}; y_G = \frac{y_A + y_B + y_C}{3}}$.
\end{itemize}

\textbf{2. Phương pháp giải các dạng}
%% DANG 0H9?1-1
\textbf{Dạng 1. Tọa độ điểm, độ dài đại số của véctơ trên 1 trục}
\begin{enumerate}
\item Gọi véctơ đơn vị trên trục là $\overrightarrow{e}$ (như $\overrightarrow{i}$ trên $Ox$, $\overrightarrow{j}$ trên $Oy$).
\item Áp dụng định nghĩa: $\overrightarrow{AB} = a\overrightarrow{e}$ thì độ dài đại số $\overline{AB} = a$.
\item Sử dụng tọa độ: $\overline{AB} = x_B - x_A$.
\end{enumerate}
Lưu ý: Độ dài đại số có thể mang dấu âm hoặc dương, khác với khoảng cách $AB = |\overline{AB}| \ge 0$.

%% DANG 0H9?1-2
\textbf{Dạng 2. Phép toán véctơ (tổng, hiệu, tích với số) trong Oxy}
\begin{enumerate}
\item Viết tọa độ các véctơ thành phần.
\item Thực hiện phép tính cộng, trừ, nhân với số theo quy tắc "hoành theo hoành, tung theo tung".
\end{enumerate}
Lưu ý: Khi tính $\overrightarrow{u} - 2\overrightarrow{v}$, nhớ nhân số $-2$ vào cả hoành và tung của $\overrightarrow{v}$ trước khi cộng để tránh nhầm dấu.

%% DANG 0H9?1-3
\textbf{Dạng 3. Tọa độ điểm và véctơ trên hệ trục Oxy}
\begin{enumerate}
\item Từ biểu diễn qua $\overrightarrow{i}, \overrightarrow{j}$, nhặt các hệ số để suy ra tọa độ véctơ/điểm.
\item Kết hợp các công thức trung điểm, trọng tâm, tọa độ véctơ để tìm điểm thỏa mãn yêu cầu.
\end{enumerate}
Lưu ý: Tính $\overrightarrow{AB}$ luôn lấy "Tọa độ điểm sau trừ tọa độ điểm trước" ($B - A$).

%% DANG 0H9?1-4
\textbf{Dạng 4. Sự cùng phương của 2 véctơ và ứng dụng}
\begin{enumerate}
\item Tính tọa độ $2$ véctơ cần xét: $\overrightarrow{u} = (x_1; y_1)$ và $\overrightarrow{v} = (x_2; y_2)$.
\item Ép điều kiện cùng phương: $x_1 y_2 - x_2 y_1 = 0$ (hoặc $\frac{x_1}{x_2} = \frac{y_1}{y_2}$ nếu các mẫu khác $0$).
\item Ứng dụng chứng minh $3$ điểm $A, B, C$ thẳng hàng $\Leftrightarrow \overrightarrow{AB}$ và $\overrightarrow{AC}$ cùng phương.
\end{enumerate}
Lưu ý: Nếu dùng tỉ số $\frac{x_1}{x_2} = \frac{y_1}{y_2}$, phải chắc chắn $x_2 \neq 0$ và $y_2 \neq 0$, nếu không hãy dùng phép nhân chéo.

%% DANG 0H9?1-5
\textbf{Dạng 5. Phân tích một véctơ theo 2 véctơ không cùng phương}
\begin{enumerate}
\item Đặt véctơ cần phân tích $\overrightarrow{c} = m\overrightarrow{a} + n\overrightarrow{b}$.
\item Lập hệ phương trình bậc nhất $2$ ẩn: $x_c = m.x_a + n.x_b$ và $y_c = m.y_a + n.y_b$.
\item Giải hệ phương trình tìm cặp số $m, n$.
\end{enumerate}
Lưu ý: Khi nhập máy tính giải hệ, chú ý thứ tự các hệ số của $m, n$ và cột hệ số tự do tương ứng.

%% DANG 0H9?1-6
\textbf{Dạng 6. Toán thực tế dùng hệ tọa độ}
\begin{enumerate}
\item Chọn hệ trục tọa độ $Oxy$ phù hợp, gắn gốc $O$ vào vị trí đặc biệt (như ngã tư, tâm mô hình).
\item Xác định tọa độ các điểm liên quan dựa vào khoảng cách thực tế.
\item Đưa bài toán về các dạng véctơ quen thuộc để tính toán.
\end{enumerate}
Lưu ý: Phải đồng nhất đơn vị đo lường (m, km,...) trên thực tế tương ứng với $1$ đơn vị độ dài trên hệ trục.

%%% BAI 0H9?2
\textbf{1. Kiến thức cần nhớ}
\begin{itemize}
\item Cho hai véctơ $\overrightarrow{a} = (x_1; y_1)$ và $\overrightarrow{b} = (x_2; y_2)$.
\item Tích vô hướng: $\boxed{\overrightarrow{a} \cdot \overrightarrow{b} = x_1 x_2 + y_1 y_2}$.
\item Độ dài véctơ: $|\overrightarrow{a}| = \sqrt{x_1^2 + y_1^2}$.
\item Khoảng cách giữa 2 điểm $A, B$: $AB = |\overrightarrow{AB}| = \sqrt{(x_B - x_A)^2 + (y_B - y_A)^2}$.
\item Góc giữa hai véctơ: $\cos(\overrightarrow{a}, \overrightarrow{b}) = \frac{\overrightarrow{a} \cdot \overrightarrow{b}}{|\overrightarrow{a}| |\overrightarrow{b}|} = \frac{x_1 x_2 + y_1 y_2}{\sqrt{x_1^2 + y_1^2} \cdot \sqrt{x_2^2 + y_2^2}}$.
\item Điều kiện vuông góc: $\overrightarrow{a} \perp \overrightarrow{b} \Leftrightarrow \overrightarrow{a} \cdot \overrightarrow{b} = 0 \Leftrightarrow \boxed{x_1 x_2 + y_1 y_2 = 0}$.
\end{itemize}

\textbf{2. Phương pháp giải các dạng}
%% DANG 0H9?2-1
\textbf{Dạng 1. Tích vô hướng, góc giữa 2 véctơ}
\begin{enumerate}
\item Xác định đúng tọa độ hai véctơ $\overrightarrow{a}$ và $\overrightarrow{b}$.
\item Thay vào công thức tích vô hướng: hoành nhân hoành cộng tung nhân tung.
\item Áp dụng công thức tính cosin góc và dùng máy tính (Shift Cos) để suy ra góc.
\end{enumerate}
Lưu ý: Kết quả của tích vô hướng là một \textbf{số thực}, không phải là một véctơ.

%% DANG 0H9?2-2
\textbf{Dạng 2. Tìm góc nhờ tích vô hướng}
\begin{enumerate}
\item Nếu bài hỏi góc của tam giác (ví dụ góc $A$), hãy tính $\overrightarrow{AB}$ và $\overrightarrow{AC}$.
\item Tính $\cos A = \cos(\overrightarrow{AB}, \overrightarrow{AC}) = \frac{\overrightarrow{AB} \cdot \overrightarrow{AC}}{AB \cdot AC}$.
\item Bấm máy tính suy ra số đo góc.
\end{enumerate}
Lưu ý: Góc giữa 2 véctơ có thể tù (nếu tích vô hướng âm). Đừng quên chung gốc khi tính góc tam giác (VD: Góc $A$ phải dùng $\overrightarrow{AB}$ và $\overrightarrow{AC}$, cấm dùng $\overrightarrow{BA}$).

%% DANG 0H9?2-3
\textbf{Dạng 3. Đẳng thức về tích vô hướng hoặc độ dài}
\begin{enumerate}
\item Chuyển các đại lượng hình học thành tọa độ.
\item Tính toán riêng từng vế bằng các công thức $\overrightarrow{a} \cdot \overrightarrow{b}$ và $|\overrightarrow{a}|$.
\item Biến đổi đại số để thu được đẳng thức cần chứng minh.
\end{enumerate}
Lưu ý: Sử dụng linh hoạt bình phương vô hướng: $\overrightarrow{a}^2 = |\overrightarrow{a}|^2 = x_1^2 + y_1^2$ để làm mất căn bậc hai.

%% DANG 0H9?2-4
\textbf{Dạng 4. Điều kiện vuông góc}
\begin{enumerate}
\item Tìm tọa độ $2$ véctơ liên quan (ví dụ $AB \perp CD \Rightarrow$ xét $\overrightarrow{AB}$ và $\overrightarrow{CD}$).
\item Lập phương trình $x_1 x_2 + y_1 y_2 = 0$.
\item Giải phương trình tìm tham số chưa biết.
\end{enumerate}
Lưu ý: Nếu bài toán yêu cầu $\Delta ABC$ vuông tại $A$, chỉ cần lập $\overrightarrow{AB} \cdot \overrightarrow{AC} = 0$. Đừng quên điều kiện các điểm phải phân biệt (véctơ khác $\overrightarrow{0}$).

%% DANG 0H9?2-5
\textbf{Dạng 5. Các bài toán tìm điểm và tập hợp điểm}
\begin{enumerate}
\item Gọi điểm cần tìm $M(x; y)$. Nếu $M \in Ox \Rightarrow M(x; 0)$, nếu $M \in Oy \Rightarrow M(0; y)$.
\item Lập phương trình dựa trên giả thiết: khoảng cách ($MA = MB$), góc, hoặc vuông góc.
\item Giải phương trình tìm tọa độ điểm hoặc biến đổi về dạng phương trình đường (đường thẳng/đường tròn).
\end{enumerate}
Lưu ý: Nếu liên quan đến khoảng cách, nên bình phương hai vế ngay từ đầu để tránh dấu căn phức tạp.

%% DANG 0H9?2-6
\textbf{Dạng 6. Cực trị và chứng minh bất đẳng thức}
\begin{enumerate}
\item Đưa biểu thức cần tìm Min/Max về một hàm số theo một ẩn $t$ (ẩn của tọa độ).
\item Biến đổi hàm số về dạng bình phương cộng hằng số $f(t) = a(t - t_0)^2 + m$ để đánh giá.
\item Có thể dùng BĐT hình học như $|\overrightarrow{u}| + |\overrightarrow{v}| \ge |\overrightarrow{u} + \overrightarrow{v}|$ để giải nhanh.
\end{enumerate}
Lưu ý: Luôn phải chỉ ra dấu "=" xảy ra khi nào (tọa độ điểm là bao nhiêu) để kết luận cực trị.

%% DANG 0H9?2-7
\textbf{Dạng 7. Toán thực tế, liên môn}
\begin{enumerate}
\item Gắn trục tọa độ $Oxy$ cho bài toán thực tế (Vật lý: Lực, vận tốc).
\item Biểu diễn các đại lượng Vật lý thành tọa độ véctơ.
\item Tính Công sinh ra $A = \overrightarrow{F} \cdot \overrightarrow{d}$ hoặc tổng hợp lực qua quy tắc tọa độ.
\end{enumerate}
Lưu ý: Công cơ học có thể âm (công cản) nếu góc tạo bởi lực và hướng dịch chuyển là góc tù, điều này hoàn toàn khớp với tính chất của tích vô hướng.

\par\medskip{\large\bfseries\color{mauPhieu} Bài 2. Tích vô hướng (theo tọa độ)}\par\smallskip
%%% BAI 0H9?1
\textbf{1. Kiến thức cần nhớ}
\begin{itemize}
\item Hệ trục tọa độ $Oxy$ có hai véctơ đơn vị $\overrightarrow{i}, \overrightarrow{j}$ tương ứng trên $Ox, Oy$.
\item Tọa độ véctơ: $\overrightarrow{u} = (x;y) \Leftrightarrow \overrightarrow{u} = x\overrightarrow{i} + y\overrightarrow{j}$.
\item Tọa độ điểm: $M(x_M; y_M) \Leftrightarrow \overrightarrow{OM} = x_M\overrightarrow{i} + y_M\overrightarrow{j}$.
\item Véctơ giữa hai điểm $A, B$: $\overrightarrow{AB} = (x_B - x_A; y_B - y_A)$.
\item Các phép toán véctơ: Cho $\overrightarrow{a} = (x_1; y_1)$ và $\overrightarrow{b} = (x_2; y_2)$.
  \begin{itemize}
  \item $\overrightarrow{a} \pm \overrightarrow{b} = (x_1 \pm x_2; y_1 \pm y_2)$.
  \item $k\overrightarrow{a} = (kx_1; ky_1)$ với $k \in \mathbb{R}$.
  \item $\overrightarrow{a} = \overrightarrow{b} \Leftrightarrow x_1 = x_2$ và $y_1 = y_2$.
  \end{itemize}
\item Hai véctơ $\overrightarrow{a}, \overrightarrow{b}$ ($\overrightarrow{b} \neq \overrightarrow{0}$) cùng phương $\Leftrightarrow \exists k \in \mathbb{R}$ sao cho $\overrightarrow{a} = k\overrightarrow{b} \Leftrightarrow x_1 = k x_2$ và $y_1 = k y_2$.
\item Tọa độ trung điểm $I$ của $AB$: $\boxed{x_I = \frac{x_A + x_B}{2}; y_I = \frac{y_A + y_B}{2}}$.
\item Tọa độ trọng tâm $G$ của $\Delta ABC$: $\boxed{x_G = \frac{x_A + x_B + x_C}{3}; y_G = \frac{y_A + y_B + y_C}{3}}$.
\end{itemize}

\textbf{2. Phương pháp giải các dạng}
%% DANG 0H9?1-1
\textbf{Dạng 1. Tọa độ điểm, độ dài đại số của véctơ trên 1 trục}
\begin{enumerate}
\item Gọi véctơ đơn vị trên trục là $\overrightarrow{e}$ (như $\overrightarrow{i}$ trên $Ox$, $\overrightarrow{j}$ trên $Oy$).
\item Áp dụng định nghĩa: $\overrightarrow{AB} = a\overrightarrow{e}$ thì độ dài đại số $\overline{AB} = a$.
\item Sử dụng tọa độ: $\overline{AB} = x_B - x_A$.
\end{enumerate}
Lưu ý: Độ dài đại số có thể mang dấu âm hoặc dương, khác với khoảng cách $AB = |\overline{AB}| \ge 0$.

%% DANG 0H9?1-2
\textbf{Dạng 2. Phép toán véctơ (tổng, hiệu, tích với số) trong Oxy}
\begin{enumerate}
\item Viết tọa độ các véctơ thành phần.
\item Thực hiện phép tính cộng, trừ, nhân với số theo quy tắc "hoành theo hoành, tung theo tung".
\end{enumerate}
Lưu ý: Khi tính $\overrightarrow{u} - 2\overrightarrow{v}$, nhớ nhân số $-2$ vào cả hoành và tung của $\overrightarrow{v}$ trước khi cộng để tránh nhầm dấu.

%% DANG 0H9?1-3
\textbf{Dạng 3. Tọa độ điểm và véctơ trên hệ trục Oxy}
\begin{enumerate}
\item Từ biểu diễn qua $\overrightarrow{i}, \overrightarrow{j}$, nhặt các hệ số để suy ra tọa độ véctơ/điểm.
\item Kết hợp các công thức trung điểm, trọng tâm, tọa độ véctơ để tìm điểm thỏa mãn yêu cầu.
\end{enumerate}
Lưu ý: Tính $\overrightarrow{AB}$ luôn lấy "Tọa độ điểm sau trừ tọa độ điểm trước" ($B - A$).

%% DANG 0H9?1-4
\textbf{Dạng 4. Sự cùng phương của 2 véctơ và ứng dụng}
\begin{enumerate}
\item Tính tọa độ $2$ véctơ cần xét: $\overrightarrow{u} = (x_1; y_1)$ và $\overrightarrow{v} = (x_2; y_2)$.
\item Ép điều kiện cùng phương: $x_1 y_2 - x_2 y_1 = 0$ (hoặc $\frac{x_1}{x_2} = \frac{y_1}{y_2}$ nếu các mẫu khác $0$).
\item Ứng dụng chứng minh $3$ điểm $A, B, C$ thẳng hàng $\Leftrightarrow \overrightarrow{AB}$ và $\overrightarrow{AC}$ cùng phương.
\end{enumerate}
Lưu ý: Nếu dùng tỉ số $\frac{x_1}{x_2} = \frac{y_1}{y_2}$, phải chắc chắn $x_2 \neq 0$ và $y_2 \neq 0$, nếu không hãy dùng phép nhân chéo.

%% DANG 0H9?1-5
\textbf{Dạng 5. Phân tích một véctơ theo 2 véctơ không cùng phương}
\begin{enumerate}
\item Đặt véctơ cần phân tích $\overrightarrow{c} = m\overrightarrow{a} + n\overrightarrow{b}$.
\item Lập hệ phương trình bậc nhất $2$ ẩn: $x_c = m.x_a + n.x_b$ và $y_c = m.y_a + n.y_b$.
\item Giải hệ phương trình tìm cặp số $m, n$.
\end{enumerate}
Lưu ý: Khi nhập máy tính giải hệ, chú ý thứ tự các hệ số của $m, n$ và cột hệ số tự do tương ứng.

%% DANG 0H9?1-6
\textbf{Dạng 6. Toán thực tế dùng hệ tọa độ}
\begin{enumerate}
\item Chọn hệ trục tọa độ $Oxy$ phù hợp, gắn gốc $O$ vào vị trí đặc biệt (như ngã tư, tâm mô hình).
\item Xác định tọa độ các điểm liên quan dựa vào khoảng cách thực tế.
\item Đưa bài toán về các dạng véctơ quen thuộc để tính toán.
\end{enumerate}
Lưu ý: Phải đồng nhất đơn vị đo lường (m, km,...) trên thực tế tương ứng với $1$ đơn vị độ dài trên hệ trục.

%%% BAI 0H9?2
\textbf{1. Kiến thức cần nhớ}
\begin{itemize}
\item Cho hai véctơ $\overrightarrow{a} = (x_1; y_1)$ và $\overrightarrow{b} = (x_2; y_2)$.
\item Tích vô hướng: $\boxed{\overrightarrow{a} \cdot \overrightarrow{b} = x_1 x_2 + y_1 y_2}$.
\item Độ dài véctơ: $|\overrightarrow{a}| = \sqrt{x_1^2 + y_1^2}$.
\item Khoảng cách giữa 2 điểm $A, B$: $AB = |\overrightarrow{AB}| = \sqrt{(x_B - x_A)^2 + (y_B - y_A)^2}$.
\item Góc giữa hai véctơ: $\cos(\overrightarrow{a}, \overrightarrow{b}) = \frac{\overrightarrow{a} \cdot \overrightarrow{b}}{|\overrightarrow{a}| |\overrightarrow{b}|} = \frac{x_1 x_2 + y_1 y_2}{\sqrt{x_1^2 + y_1^2} \cdot \sqrt{x_2^2 + y_2^2}}$.
\item Điều kiện vuông góc: $\overrightarrow{a} \perp \overrightarrow{b} \Leftrightarrow \overrightarrow{a} \cdot \overrightarrow{b} = 0 \Leftrightarrow \boxed{x_1 x_2 + y_1 y_2 = 0}$.
\end{itemize}

\textbf{2. Phương pháp giải các dạng}
%% DANG 0H9?2-1
\textbf{Dạng 1. Tích vô hướng, góc giữa 2 véctơ}
\begin{enumerate}
\item Xác định đúng tọa độ hai véctơ $\overrightarrow{a}$ và $\overrightarrow{b}$.
\item Thay vào công thức tích vô hướng: hoành nhân hoành cộng tung nhân tung.
\item Áp dụng công thức tính cosin góc và dùng máy tính (Shift Cos) để suy ra góc.
\end{enumerate}
Lưu ý: Kết quả của tích vô hướng là một \textbf{số thực}, không phải là một véctơ.

%% DANG 0H9?2-2
\textbf{Dạng 2. Tìm góc nhờ tích vô hướng}
\begin{enumerate}
\item Nếu bài hỏi góc của tam giác (ví dụ góc $A$), hãy tính $\overrightarrow{AB}$ và $\overrightarrow{AC}$.
\item Tính $\cos A = \cos(\overrightarrow{AB}, \overrightarrow{AC}) = \frac{\overrightarrow{AB} \cdot \overrightarrow{AC}}{AB \cdot AC}$.
\item Bấm máy tính suy ra số đo góc.
\end{enumerate}
Lưu ý: Góc giữa 2 véctơ có thể tù (nếu tích vô hướng âm). Đừng quên chung gốc khi tính góc tam giác (VD: Góc $A$ phải dùng $\overrightarrow{AB}$ và $\overrightarrow{AC}$, cấm dùng $\overrightarrow{BA}$).

%% DANG 0H9?2-3
\textbf{Dạng 3. Đẳng thức về tích vô hướng hoặc độ dài}
\begin{enumerate}
\item Chuyển các đại lượng hình học thành tọa độ.
\item Tính toán riêng từng vế bằng các công thức $\overrightarrow{a} \cdot \overrightarrow{b}$ và $|\overrightarrow{a}|$.
\item Biến đổi đại số để thu được đẳng thức cần chứng minh.
\end{enumerate}
Lưu ý: Sử dụng linh hoạt bình phương vô hướng: $\overrightarrow{a}^2 = |\overrightarrow{a}|^2 = x_1^2 + y_1^2$ để làm mất căn bậc hai.

%% DANG 0H9?2-4
\textbf{Dạng 4. Điều kiện vuông góc}
\begin{enumerate}
\item Tìm tọa độ $2$ véctơ liên quan (ví dụ $AB \perp CD \Rightarrow$ xét $\overrightarrow{AB}$ và $\overrightarrow{CD}$).
\item Lập phương trình $x_1 x_2 + y_1 y_2 = 0$.
\item Giải phương trình tìm tham số chưa biết.
\end{enumerate}
Lưu ý: Nếu bài toán yêu cầu $\Delta ABC$ vuông tại $A$, chỉ cần lập $\overrightarrow{AB} \cdot \overrightarrow{AC} = 0$. Đừng quên điều kiện các điểm phải phân biệt (véctơ khác $\overrightarrow{0}$).

%% DANG 0H9?2-5
\textbf{Dạng 5. Các bài toán tìm điểm và tập hợp điểm}
\begin{enumerate}
\item Gọi điểm cần tìm $M(x; y)$. Nếu $M \in Ox \Rightarrow M(x; 0)$, nếu $M \in Oy \Rightarrow M(0; y)$.
\item Lập phương trình dựa trên giả thiết: khoảng cách ($MA = MB$), góc, hoặc vuông góc.
\item Giải phương trình tìm tọa độ điểm hoặc biến đổi về dạng phương trình đường (đường thẳng/đường tròn).
\end{enumerate}
Lưu ý: Nếu liên quan đến khoảng cách, nên bình phương hai vế ngay từ đầu để tránh dấu căn phức tạp.

%% DANG 0H9?2-6
\textbf{Dạng 6. Cực trị và chứng minh bất đẳng thức}
\begin{enumerate}
\item Đưa biểu thức cần tìm Min/Max về một hàm số theo một ẩn $t$ (ẩn của tọa độ).
\item Biến đổi hàm số về dạng bình phương cộng hằng số $f(t) = a(t - t_0)^2 + m$ để đánh giá.
\item Có thể dùng BĐT hình học như $|\overrightarrow{u}| + |\overrightarrow{v}| \ge |\overrightarrow{u} + \overrightarrow{v}|$ để giải nhanh.
\end{enumerate}
Lưu ý: Luôn phải chỉ ra dấu "=" xảy ra khi nào (tọa độ điểm là bao nhiêu) để kết luận cực trị.

%% DANG 0H9?2-7
\textbf{Dạng 7. Toán thực tế, liên môn}
\begin{enumerate}
\item Gắn trục tọa độ $Oxy$ cho bài toán thực tế (Vật lý: Lực, vận tốc).
\item Biểu diễn các đại lượng Vật lý thành tọa độ véctơ.
\item Tính Công sinh ra $A = \overrightarrow{F} \cdot \overrightarrow{d}$ hoặc tổng hợp lực qua quy tắc tọa độ.
\end{enumerate}
Lưu ý: Công cơ học có thể âm (công cản) nếu góc tạo bởi lực và hướng dịch chuyển là góc tù, điều này hoàn toàn khớp với tính chất của tích vô hướng.
